Market infrastructures can be grouped in two categories:
dealer markets (also known as quote-driven markets)
and
limit order markets (also known as order-driven markets)
-- see \citealp[Chapter 1]{FPR13mar}.
Several trading venues around the world have adopted infrastructures that fall in the latter category.
\href{https://www.euronext.com/}{EURONEXT}\footnote{See
\url{https://www.euronext.com/}.},
the
\href{https://www.six-group.com/}{Swiss stock exchange}\footnote{See
\url{https://www.six-group.com/}.},
\href{http://www.nasdaqomxnordic.com/}{Helsinki stock exchange}\footnote{See
\url{http://www.nasdaqomxnordic.com/}.},
the
\href{https://www.thecse.com/}{Canadian stock exchange}\footnote{See
\url{https://www.thecse.com/}.}
and
\href{https://www.tsx.com/}{Toronto stock exchange}\footnote{See
\url{https://www.tsx.com/}.},
the
\href{https://www2.asx.com.au/}{Australian Securities Exchange}\footnote{See
\url{https://www2.asx.com.au/}.},
and the stock exchanges in
\href{https://www.hkex.com.hk/}{Hong Kong}\footnote{See
\url{https://www.hkex.com.hk/}.},
\href{http://www.szse.cn/}{Shenzhen}\footnote{See
\url{http://www.szse.cn/}.}
and
\href{https://www.jpx.co.jp/}{Tokyo}\footnote{See
\url{https://www.jpx.co.jp/}.},
all of them operate as purely order-driven markets;
\href{https://www.nyse.com/}{NYSE}\footnote{See
\url{https://www.nyse.com/}.},
\href{https://www.nasdaq.com/}{NASDAQ}\footnote{See
\url{https://www.nasdaq.com/}.}
and
\href{https://www.londonstockexchange.com/}{LSE}\footnote{See
\url{https://www.londonstockexchange.com/}.}
have order-driven market infrastructure
with some legacy quote-driven mechanisms.

Order-driven markets are transparent.
Indeed,
all agents participating in an order-driven market observe in real-time all other participants'
commitments to trade,
which are organised in the so-called \emph{limit order book}.

Limit order books provide high-frequency financial data that are susceptible of various
quantitative investigations.
A survey of empirical and theoretical studies on limit order books is \cite{GPWMFH13lim}.

From a stochastic modelling perspective, important features of these datasets include:
\emph{i.} the continuous physical time of the events recorded;
\emph{ii.} the complex relations that may emerge from the occurrence in time
of one type of event and the time distribution of other types of events.

Order-driven markets are trading venues organised around \emph{limit orders}.
A limit order is the fundamental action that market participants in order-driven markets can perform.
It is represented by a 4-tuple $(t,q,p,d)$,
where $t$ denotes time, $q$ denotes size, $p$ denotes price, and $d$ denotes direction.
We say that a market participant \emph{posts} a limit order $(t,q,p,d)$
if at time $t$ she submits to the exchange her commitment to buy ($d=1$) or sell ($d=-1$)
the amount $q$ at the price $p$.
The price $p$ is interpreted as the highest price at which she is committed to buy if $d=1$,
or the lowest price at which she is committed to sell if $d=-1$.
By regulation, $p$ must be an integer multiple of some fixed $\tickSizeOfLOB>0$.
Market participants have also the possibility to withdraw their commitments to trade,
by \emph{cancelling} a previously submitted limit order.

A trading epoch in an order-driven market is the collection of all the limit orders
submitted by market participants within some time interval.
Limit orders cannot be submitted simultaneously,
so that a limit order is identified by its timestamp $t$.
We express this mathematically by saying that a trading epoch is the graph of a function
from a subset of the positive half-line to the space
$\lbrace (q,p,d): \, q\geq0, p\geq0, d=-1,1\rbrace$.

When a market participant submits a limit order $(t,q,p,d)$,
a matching algorithm run by the trading platform looks for possible counterparts to her order:
if other participants' commitments to trade exist to (partially) fulfill her order
(at a price not worse than $p$ from her point of view),
then her order is (partially) cleared.
The fraction of her order that is cleared disappears from the market,
and it is said to be \emph{executed};
the fraction that could not be fulfilled is recorded in the \emph{limit order book},
waiting for counterparts to trade with.

A limit order $(t,q,p,d)$ in a trading epoch is said to be \emph{active} (or outstanding) at time $u$
if $t\leq u$ and by time $u$ the order has neither been executed nor been cancelled.

The search for possible trading counterparts for the matching of the incoming order $(t,q,p,d)$
happens among the active orders with timestamp $s<t$ and opposite direction $-d$.
We say that the active order $(s,\rho,\pi,-d)$ matches the incoming order $(t,q,p,d)$
if
$\pi d \leq p d$;
in this case, the two orders are executed one against the other,
and $\rho \wedge q$ units of shares are transacted at the price $\pi$ per unit of share.
After the match, the matched orders are replaced with the orders
$(s,\tilde{\rho}, \pi, -d)$ and $(t,\tilde{q}, p, d)$,
where $\tilde{\rho} = \rho - \rho \wedge q$
and $\tilde{q} = q - \rho \wedge q$.
At least one of them has null size.
If $\tilde{\rho} = 0$, i.e. $q \geq \rho$,
the order with time stamp $s$ ceases to be active and it is deleted from its queue.
If $\tilde{q} = 0$, i.e. $q\leq \rho$,
then the order with time stamp $t$ (i.e. the incoming order) is fully executed,
and the search for possible matching counterparts stops.
If instead $\tilde{q} > 0$, i.e. $q>\rho$,
then the search continues among active orders with opposite direction $-d$.

The search follows the priority given by the ordering
\begin{equation}
\label{eq.orderingOfLimitOrders}
 (s,\rho,\pi,-d)<(s\derivative,\rho\derivative,\pi\derivative,-d)
 \qquad
 \text{ if and only if }
 \qquad
 (d\pi<d\pi\derivative)\vee((\pi=\pi\derivative)\wedge(s<s\derivative)),
\end{equation}
whereby $(s,\rho,\pi,-d)$ is given priority over $(s\derivative,\rho\derivative,\pi\derivative,-d)$
if $ d\pi<d\pi\derivative$ or, in case $\pi=\pi\derivative$, $s<s\derivative$.
This means that if both $(s,\rho,\pi,-d)$ and $(s\derivative,\rho\derivative,\pi\derivative,-d)$
are potential matches to the incoming order $(t, q,p,d)$
(i.e.
$\pi d \leq p d$ and
$\pi\derivative d \leq p d$),
then
$(t, q, p, d)$ is executed first against $(s, \rho, \pi, -d)$
and only its outstanding part $(t, \tilde{q}, p, d)$ after this transaction
is executed against $(s\derivative, \rho\derivative, \pi\derivative, -d)$.

\begin{lemma}
\label{lemma.orderingOfOrders}
Let $\tradingEpoch$ be a trading epoch in an order-driven market.
Then, equation \eqref{eq.orderingOfLimitOrders} defines a total order in
$\tradingEpoch\cap \lbrace (s,\rho,\pi,-1): s\geq0, \rho\geq0,\pi\geq0\rbrace$,
and it defines a total order in
$\tradingEpoch\cap \lbrace (s,\rho,\pi,1): s\geq0, \rho\geq0,\pi\geq0\rbrace$.
\end{lemma}
\begin{proof}
 Antisymmetry and transitivity follow from \eqref{eq.orderingOfLimitOrders} alone,
 being those of the lexicographic order on the pair $(d\pi, s)$.
 Connexity is where the trading epoch enters.
 Let two distinct orders of the same side be given.
 If their prices differ, then $d\pi < d\pi\derivative$ or $d\pi\derivative < d\pi$,
 and \eqref{eq.orderingOfLimitOrders} decides between them.
 If their prices agree, their timestamps must differ,
 since $\tradingEpoch$ is the graph of a function defined on a subset of $\R$
 and therefore contains no two orders with the same timestamp;
 and \eqref{eq.orderingOfLimitOrders} then decides between them by $s$.
\end{proof}
Lemma \ref{lemma.orderingOfOrders} justifies the following definition.
\begin{defi}
\label{def.orderQueue}
 Let $\tradingEpoch$ be a trading epoch in an order-driven market.
 The ask order queue $\askOrderQueue_t$ at time $t$ is defined as
 \begin{equation*}
  \askOrderQueue_t := \lbrace (s,q,p,-1)\in \tradingEpoch : \, (s,q,p,-1) \text{ is active at time } t \rbrace.
 \end{equation*}
 Similarly, the bid order queue $\bidOrderQueue_t$ at time $t$ is defined as
 \begin{equation*}
  \bidOrderQueue_t := \lbrace (s,q,p,+1) \in \tradingEpoch : \, (s,q,p,+1) \text{ is active at time } t \rbrace.
 \end{equation*}
\end{defi}

A limit order book is a grid of equally spaced prices at which active limit orders sit.
The space between consecutive grid nodes is called the tick size of the LOB,
which we denote by $\tickSizeOfLOB$.
As mentioned above,
all orders submitted to the exchange must have prices that are integer multiples of the tick size.
Prices are increasing from left to right.
At every node of the grid,
outstanding limit orders to buy or sell at the corresponding price are collected;
these limit orders are also said to be queuing there.
Spontaneously -- i.e.
by market forces -- such orders organise in a way that
buy offers will be displaced on the left (the so-called \emph{bid side})
and sell offers will be on the right (the so-called \emph{ask side}).

Indeed,
if there were a buy (respectively, sell) limit order to the right (left) of some sell (buy) limit orders,
or at the same node,
then the matching algorithm would have matched them, clearing them out from the order book.

Therefore, the whole configuration of a limit order book at time $t$ is given
when the following variables are specified:
\begin{enumerate}[label={\emph{\roman{*}}.}, ref={\emph{\roman{*}}.}]
\item\label{item.bestAskPrice} the best ask price $\bestAskPrice_{t}$, i.e.
the lowest price at which one can find sell offers active at time $t$;
\item\label{item.bestBidPrice} the best bid price $\bestBidPrice_t$, i.e.
the highest price at which one can find buy offers active at time $t$;
\item\label{item.askSize} the size $\nthBestAskSize[i]_t$ of ask offers at price
$\nthBestAskPrice[i]_t= \bestAskPrice_t + (i-1)\tickSizeOfLOB$, for $i=1,2,\dots$, i.e.
the quantity
\begin{equation*}
 \nthBestAskSize[i]_t = \sum_{\lbrace(s,q,\nthBestAskPrice[i]_t,-1) \in \askOrderQueue_t : \quad s\leq t\rbrace} q,
\end{equation*}
where the sum is over all the active sell limit orders submitted by time $t$
for a price of $\bestAskPrice_t + (i-1)\tickSizeOfLOB$;
\item\label{item.bidSize} the size $\nthBestBidSize[i]_t$ of bid offers at price
$\nthBestBidPrice[i]_t=\bestBidPrice_t - (i-1)\tickSizeOfLOB$, for $i=1,2,\dots$, i.e.
the quantity
\begin{equation*}
 \nthBestBidSize[i]_t = \sum_{\lbrace (s,q,\nthBestBidPrice[i]_t,1) \in \bidOrderQueue_t: \quad s\leq t\rbrace} q,
\end{equation*}
where the sum is over all the active buy limit orders submitted by time $t$
for a price of $\bestBidPrice_t - (i-1)\tickSizeOfLOB$.
\end{enumerate}
Points \ref{item.bestAskPrice}-\ref{item.bidSize} above describe how the configuration at time $t$
of the limit order book is reconstructed from the ask queue $\askOrderQueue_t$
and the bid queue $\bidOrderQueue_t$.


\begin{remark}
\label{remark.sizeVersusVolume}
In points \ref{item.askSize} and \ref{item.bidSize} the sum of the quantities queuing at a
price level is called the \emph{size} of that level,
measured in number of shares.
We reserve \emph{volume} for the number of shares transacted over an interval:
a level has a size at every instant, whereas volume accumulates between two instants.
Volume is not a change in size.
A level shrinks by cancellation as well as by execution,
so no sequence of sizes determines the volume traded;
by Proposition \ref{prop.decompositionOfLimitOrder},
the volume over a window is the sum of $\marketOrderSize$
over the market orders submitted within it.
\end{remark}

The whole configuration of a limit order book at time $t$ is described by the variables
$(\bestAskPrice_t, \bestBidPrice_t, \lbrace (\nthBestAskSize[i]_t, \nthBestBidSize[i]_t) : \, \, i=1,2,\dots\rbrace)$.
Therefore,
in this course
limit order books will always mean aggregated order book.
This means that we will not map the queue sizes to  the orders in the trading epoch.


We will use derived quantities from this configuration:
the bid-ask spread,
the mid-price,
the queue-imbalance,
the micro-price,
and the sweep cost.


\begin{defi}
The spread at time $t$ is the distance $\abs{\bestAskPrice_t - \bestBidPrice_t}$
between best ask price and best bid price,
which we denote by $\LOBspread_t$.
\end{defi}
Since the two sides never overlap and every price is an integer multiple of $\tickSizeOfLOB$,
the spread satisfies $\LOBspread_t \geq \tickSizeOfLOB$ whenever both sides are non-empty.
\begin{defi}
The mid-price at time $t$ is the mid point in between $\bestBidPrice_t$ and $\bestAskPrice_{t}$,
which we denote by $\midPrice_t$ and is given by
$\midPrice_t = (\bestAskPrice_t + \bestBidPrice_{t})/{2}$.
\end{defi}
\begin{defi}
The $n$-levels queue imbalance at time $t$,
denoted $\queueImbalance[n]_t$,
is the normalised excess of limit orders on the first $n$ levels of the bid side
compared to the limit orders on the first $n$ levels of the ask side, namely
\begin{equation}
\label{eq.queueImbalance}
\queueImbalance[n]_t =
\frac{\sum_{i\leq n}\nthBestBidSize[i]_t - \sum_{i\leq n}\nthBestAskSize[i]_t}{\sum_{i\leq n}\nthBestBidSize[i]_t + \sum_{i\leq n}\nthBestAskSize[i]_t},
\end{equation}
where $\sum_{i\leq n}\nthBestBidSize[i]_t$
(respectively, $\sum_{i\leq n}\nthBestAskSize[i]_t$)
is the cumulative size on the first $n$ bid (ask) levels.
\end{defi}
\begin{defi}
The micro-price at time $t$ is
\begin{equation}
\label{eq.microPrice}
\microPrice_t
=
\midPrice_t + \frac{\LOBspread_t}{2}\queueImbalance[1]_t .
\end{equation}
\end{defi}
\begin{remark}
Notice that \eqref{eq.microPrice} is the weighting of the two best prices
in which each price carries the size resting on the \emph{opposite} side,
\begin{equation*}
\microPrice_t
=
\frac{\nthBestBidSize[1]_t\bestAskPrice_t + \nthBestAskSize[1]_t\bestBidPrice_t}
{\nthBestBidSize[1]_t + \nthBestAskSize[1]_t} ,
\end{equation*}
so that $\microPrice_t$ lies in $[\bestBidPrice_t, \bestAskPrice_t]$
and a bid-heavy book carries it towards the ask.\footnote{
The name is also given, in \cite{Sto18mic}, to a different object:
the martingale-adjusted limit of the mid-price conditional on the imbalance.
The two are not the same, and the literature does not always keep them apart.}
\end{remark}
\begin{defi}
Let the ask side hold at least $\quantityToLiquidate > 0$ shares at time $t$,
and for $i=1,2,\dots$ let
\begin{equation}
\label{eq.fillAtLevel}
q_i
:=
\Big( \quantityToLiquidate - \sum_{j<i} q_j \Big)_{+} \wedge \nthBestAskSize[i]_t
\end{equation}
be the quantity that a buy market order for $\quantityToLiquidate$ shares fills
at the $i$-th ask level, so that $\sum_{i\geq1} q_i = \quantityToLiquidate$.
The sweep cost is its per-share cost measured against the mid-price,
\begin{equation}
\label{eq.sweepCost}
\sweepCost_t(\quantityToLiquidate)
=
\frac{1}{\quantityToLiquidate}\sum_{i\geq1} q_i \big(\nthBestAskPrice[i]_t - \midPrice_t\big)
=
\frac{\LOBspread_t}{2}
+
\frac{1}{\quantityToLiquidate}\sum_{i\geq1} q_i (i - 1)\tickSizeOfLOB .
\end{equation}
For a sell market order the definition is mutatis mutandis the same,
with the bid side and $\midPrice_t - \nthBestBidPrice[i]_t$ in place of
$\nthBestAskPrice[i]_t - \midPrice_t$.
\end{defi}
The index $i$ in \eqref{eq.fillAtLevel} and \eqref{eq.sweepCost} runs over the levels of the
price grid, not over the occupied ones,
so that $\nthBestAskPrice[i]_t = \bestAskPrice_t + (i-1)\tickSizeOfLOB$
and the second equality in \eqref{eq.sweepCost} follows from
$\bestAskPrice_t - \midPrice_t = \LOBspread_t/2$.

When a limit order comes to the market,
it modifies the previous configuration of the order book.
Proposition \ref{prop.lobUpdate} gives the updating rule for the variables
$(\bestAskPrice_t, \bestBidPrice_t, \lbrace (\nthBestAskSize[i]_t, \nthBestBidSize[i]_t) : \, \, i=1,2,\dots\rbrace)$.
\begin{prop}
\label{prop.lobUpdate}
 Let $(t,q,p,-1)$ be a sell limit order that arrives to the market at time $t$.
 Let $(\bestAskPrice_{t-}, \bestBidPrice_{t-},\lbrace (\nthBestAskSize[i]_{t-}, \nthBestBidSize[i]_{t-}) : \, \, i=1,2,\dots\rbrace )$
 represent the limit order book immediately before $t$,
 and assume that $(t,q,p,-1)$ does not consume the whole of the price-eligible side,
 that is, $q < \sum_{i\geq1}\nthBestBidSize[i]_{t-}\indicator{\nthBestBidPrice[i]_{t-}\geq p}$.
 Adopt the convention $\nthBestAskSize[j]_{t-} = \nthBestBidSize[j]_{t-} = 0$ for $j \leq 0$.
 Then, prices and sizes are updated according to the following formulae:
 \begin{equation}
 \label{eq.lobUpdateSellOrder}
  \begin{split}
   \bestBidPrice_t =& \nthBestBidPrice[1+N]_{t-};
   \\
   \nthBestBidSize[k]_{t}
   =& \nthBestBidSize[k+N]_{t-}
   - \left( q^{k+N-1}-q^{k+N}\right), \qquad k=1,2,\dots ;
   \\
   \bestAskPrice_t = & \bestAskPrice_{t-}
   - \max\left(0,\bestAskPrice_{t-} - p\right) \indicator{q^{\infty} > 0} ;
   \\
   \nthBestAskSize[k]_t = &  \nthBestAskSize[k+\delta\bestAskPrice_t/\tickSizeOfLOB]_{t-}
   +q^{\infty}\indicator{p=\bestAskPrice_t + (k-1)\tickSizeOfLOB}, \qquad k=1,2,\dots ;
  \end{split}
 \end{equation}
 where $\delta\bestAskPrice_t = \bestAskPrice_t - \bestAskPrice_{t-}$
 denotes the (non-positive) jump of the best ask price;
 $N:=\max(0,N_s \wedge N_p -1)$,
 where $N_s := \inf\lbrace n\geq 0 : \, q<\sum_{i=1}^{n} \nthBestBidSize[i]_{t-}\rbrace$
 and $N_p:=\inf\lbrace n\geq 1: \, \nthBestBidPrice[n]_{t-}< p \rbrace$; and
 \begin{equation*}
  \begin{split}
   q^{i}:=& \max\left( 0,
    q-\sum_{k=1}^{i\wedge N_s \wedge (N_p - 1)} \nthBestBidSize[k]_{t-}
   \right)
   \\
   =& \max\left( 0,
   q-\sum_{1 \leq k\leq i} \nthBestBidSize[k]_{t-}\indicator{\nthBestBidPrice[k]_{t-} \geq p}
   \right)
  \end{split}
 \end{equation*}
 is the outstanding quantity to be executed after the first $i$ levels have been consumed.
 For a buy limit order $(t,q,p,+1)$, the updating rules are similar, with opposite side and direction.
\end{prop}
\begin{proof}
 We first reduce the order to a stream of elementary ones.
 By Lemma \ref{lemma.orderingOfOrders} the matching algorithm consumes the bid queue in the
 order \eqref{eq.orderingOfLimitOrders},
 and by the hypothesis it stops before exhausting the price-eligible side.
 We claim that processing $(t,q,p,-1)$ is equivalent to processing the stream
 \begin{equation}\label{eq.elementaryStream}
  (t,\nthBestBidSize[1]_{t-},\nthBestBidPrice[1]_{t-},-1), \dots,
  (t,\nthBestBidSize[N]_{t-},\nthBestBidPrice[N]_{t-},-1),
  (t,q^{N}, p, -1),
 \end{equation}
 with priority from left to right.
 The claim is proved by induction on the number of levels the order reaches.
 If it reaches none, then $N=0$, the stream is the single order $(t,q,p,-1)$ itself
 and there is nothing to prove.
 If it reaches $j \geq 1$ levels, then by \eqref{eq.orderingOfLimitOrders} the first
 $\nthBestBidSize[1]_{t-}$ shares are matched against the whole of the first bid level,
 at the price $\nthBestBidPrice[1]_{t-}$,
 before any other resting order is considered;
 this is the effect of the first element of \eqref{eq.elementaryStream},
 and what remains to be processed is the sell order
 $(t, q - \nthBestBidSize[1]_{t-}, p, -1)$
 against the book with its first bid level removed,
 which reaches $j-1$ levels and to which the induction hypothesis applies.
 Each element of \eqref{eq.elementaryStream} is therefore \emph{elementary}:
 it either clears exactly one level or rests without matching.
 It remains to check the update formulae on one such element,
 which is what the following three cases do.
 \begin{enumerate}[label={\emph{\roman{*}}.}]
  \item No match with any buy order.
  This happens if $p>\bestBidPrice_{t-}$.
  In this case therefore,
  $\bestBidPrice_t = \bestBidPrice_{t-}$ and $\nthBestBidSize[k]_{t} = \nthBestBidSize[k]_{t-}$
  are the correct update rules on the bid side, because no modification happens.
  With $N_p=1$ (so that $N=0$) and $q^i=q$ for all $i\geq 0$
  these updates are retrieved from the formulae in the statement.
  As for the ask side,
  $\bestAskPrice_t$ must be updated to the minimum between $\bestAskPrice_{t-}$ and $p$,
  and the stated updating rule yields this.
  Finally, we have that $-\delta\bestAskPrice_t/\tickSizeOfLOB$ is the number of levels
  by which the ask price is improved by the incoming order at time $t$.
  Hence, if indeed there is an improvement (i.e. $p<\bestAskPrice_{t-}$),
  we have that $\lbrace \nthBestAskSize[k]_{t}: \, k=1,2,\dots\rbrace$
  represents the sizes of $\lbrace \nthBestAskSize[k]_{t-}: \, k=1,2,\dots\rbrace$
  where the index $k$ is shifted by $-\delta\bestAskPrice_t/\tickSizeOfLOB$,
  and $\nthBestAskSize[1]_t = q$, the size of the incoming order.
  If instead there is no improvement (i.e. $p\geq \bestAskPrice_{t-}$),
  the only size $\nthBestAskSize[k]_{t-}$ that is modified
  is the one corresponding to the price $p=\bestAskPrice_t + (k-1)\tickSizeOfLOB$,
  and such modification is the addition of $q$.
  \item Complete execution at the first bid level.
  This happens if $p\leq\bestBidPrice_{t-}$ and $q<\nthBestBidSize[1]_{t-}$.
  In this case the ask side is not altered,
  and the stated updating rules yield that because $q^{\infty}=0$.
  As for the bid side, we must have
  $\bestBidPrice_t=\bestBidPrice_{t-}$,
  $\nthBestBidSize[1]_t = \nthBestBidSize[1]_{t-}- q$,
  and $\nthBestBidSize[k]_t = \nthBestBidSize[k]_{t-}$ for $k\geq 2$.
  Since in this case $N_s=1$ and $N_p\geq 2 $, $q^0 = q$ and $q^i = 0$ for all $i\geq 1$
  the stated updating rules encompass this case as well.
  \item Partial execution at first bid level.
  This happens if $p=\bestBidPrice_{t-}$ and $q\geq \nthBestBidSize[1]_{t-}$.
  In this case, we must check that the stated updating rule yield
  $\bestBidPrice_t = \nthBestBidPrice[2]_{t-}$,
  $\nthBestBidSize[k]_t = \nthBestBidSize[k+1]_{t-}$ for all $k\geq 1$, on the bid side.
  As $N_s\geq 2$, $N_p=2$, $q^0 = q$, and $q^i = q-\nthBestBidSize[1]_{t-}$ for all $i\geq 1$,
  this is indeed checked.
  Finally, on the ask side if $q=\nthBestBidSize[1]_{t-}$ there are not modifications at all.
  If instead $q>\nthBestBidSize[1]_{t-}$ the updating rules must yield
  $\bestAskPrice_t = p$, $\nthBestAskSize[1]_t = q-\nthBestBidSize[1]_{t-}$,
  and the size index $k$ shifted by the number of levels in the spread at time $t-$.
  Since with $q>\nthBestBidSize[1]_{t-}$ it holds $q^{\infty} = q-\nthBestBidSize[1]_{t-}$
  and $\delta\bestAskPrice_t = -\LOBspread_{t-}$ this is confirmed.
 \end{enumerate}
\end{proof}

The arrival of a limit order to the market triggers two events:
one is the consumption of the liquidity capable to (partially) service the limit order,
the other is the addition of the non-executed part of the limit order
to the appropriate queue in the order book.
Proposition \ref{prop.decompositionOfLimitOrder} states the terms of this decomposition.

\begin{prop}
\label{prop.decompositionOfLimitOrder}
 Given the configuration
 $(\bestAskPrice_{t-}, \bestBidPrice_{t-},\lbrace (\nthBestAskSize[i]_{t-}, \nthBestBidSize[i]_{t-}) : \, \, i=1,2,\dots\rbrace )$
 of the limit order book immediately before time $t$,
 processing the sell limit order $(t,q,p,-1)$ is equivalent to processing the pair of orders
 $[(t,\marketOrderSize,0,-1),(t,q-\marketOrderSize,p,-1) ]$,
 with the former having priority over the latter, and where
 \begin{equation*}
  \marketOrderSize:= \min\left(q, \sum_{i\geq 1} \nthBestBidSize[i]_{t-} \indicator{\nthBestBidPrice[i]_{t-}\geq p}\right).
 \end{equation*}
Similarly, processing the buy limit order $(t,q,p,+1)$ is equivalent to processing the pair of orders
$[(t,\marketOrderSize,\infty,+1),(t,q-\marketOrderSize,p,+1) ]$, where
 \begin{equation*}
  \marketOrderSize:= \min\left (q, \sum_{i\geq 1} \nthBestAskSize[i]_{t-} \indicator{\nthBestAskPrice[i]_{t-}\leq  p}\right).
 \end{equation*}
\end{prop}
\begin{proof}
 Let $(t,q,p,-1)$ be a sell limit order.
 Let $N_s$, $N_p$ and $q^i$ be as in Proposition \ref{prop.lobUpdate}.
 Let $N_s^M$, $N_p^M$ and $\marketOrderSize[i]$ be the corresponding quantities
 for the order $(t,\marketOrderSize,0,-1)$.
 Notice that $N_s^M = N_s^M \wedge N_p^M$.
 Processing $(t,\marketOrderSize,0,-1)$ does not affect the ask side
 because $\marketOrderSize[\infty] = 0$;
 moreover the effects on the bid side of processing $(t,\marketOrderSize,0,-1)$
 are the same as those of $(t,q,p,-1)$,
 because $\inf \lbrace n\geq 0 : \, \marketOrderSize < \sum_{i=1}^{n} \nthBestBidSize[i]_{t-} \rbrace = N_s \wedge N_p$
 and $q^{k+N-1} - q^{k+N} = \marketOrderSize[k+N_s^M - 1] - \marketOrderSize[k+N_s^M]$.
 Therefore, after $(t,\marketOrderSize,0,-1)$ has been processed,
 the bid side is the same as the bid side after the processing of $(t,q,p,-1)$.

 Processing $(t,q-\marketOrderSize,p,-1)$ after $(t,\marketOrderSize,0,-1)$
 does not alter the bid side because either $\bestBidPrice_{t-} < p$ or $q-\marketOrderSize=0$.
 Moreover, $(t,q-\marketOrderSize,p,-1)$ has the same effects on the ask side
 as those of $(t,q,p,-1)$ because $q-\marketOrderSize = q^{\infty}$.

 The proof of the decomposition of a buy limit order is mutatis mutandis the same.
\end{proof}

The first component $(t,\marketOrderSize,0,-1)$
in the decomposition $[(t,\marketOrderSize,0,-1),(t,q-\marketOrderSize,p,-1) ]$
of the sell limit order $(t,q,p,-1)$ is called \emph{sell market order}.
Notice that in the 4-tuple $(t,\marketOrderSize,0,-1)$, the price $p$ is set to zero,
and this guarantees immediate execution:
the sale of the amount $\marketOrderSize$ is instantaneously matched
with outstanding buy limit orders on the bid side,
and no fraction of $\marketOrderSize$ is put into the queue.
On the contrary,
the second component $(t,q-\marketOrderSize,p,-1)$ specifies exactly that part of $(t,q,p,-1)$
which will be queued.
Similarly, $(t,\marketOrderSize,\infty,+1)$ represents the market order component of
the buy limit order $(t,q,p,+1)$,
and it is referred to as \emph{buy market order}.
The price $p=\infty$ is a way to express the fact that a counterpart
for the purchase of the amount $\marketOrderSize$
will instantaneously be found on the ask side of the order book.
In the following, the term 'market order' (either buy or sell) will be referring
to the first component in the decomposition of a limit order
with non-zero market order size $\marketOrderSize >0$.\footnote{
In some trading venues, traders can actually submit market orders,
i.e. buy or sell orders that are executed without price constraints,
at least as long as offers with the opposite direction exist.
Even in such cases, we keep our convention of referring to the fraction of a limit order
that is executed upon submission as market order.}

A sell market order $(t,\marketOrderSize,0,-1)$ is said to 'walk the book'
if $\marketOrderSize > \nthBestBidSize[1]_{t-}$;
this implies $N\geq 1 $ in the notation of Proposition \ref{prop.lobUpdate},
the converse failing when the order exactly clears the best level.
Similarly, a buy market order $(t,\marketOrderSize,\infty,1)$ walks the book
if $\marketOrderSize > \nthBestAskSize[1]_{t-}$.
A market order walks the book if an only if
its sweep cost is larget than the bid-ask sprad.

Given a time window $\timeWindow$,
the evolution in time of the limit order book
$\lbrace (\bestAskPrice_t, \bestBidPrice_t,\lbrace (\nthBestAskSize[i]_t, \nthBestBidSize[i]_t) : \, \, i=1,2,\dots\rbrace): \quad 0\leq t \leq \timeHorizon \rbrace$
results from the history of all limit order submissions,
cancellations and executions that happened in $\timeWindow$.
If every limit order in $\lbrace (t,q,p,d): \, 0\leq t\leq \timeHorizon\rbrace$
is decomposed as per in Proposition \ref{prop.decompositionOfLimitOrder},
then the seller-initiated trades that happened in $\timeWindow$ are
$\lbrace (t,\marketOrderSize,0,-1):\, \, \marketOrderSize>0, \, 0\leq t\leq \timeHorizon \rbrace$,
and the buyer-initiated trades that happened in $\timeWindow$ are
$\lbrace (t,\marketOrderSize,\infty,+1):\, \, \marketOrderSize>0, \, 0\leq t\leq \timeHorizon \rbrace$.
Hence, in the following we will identify trades with market orders
of non-zero size $\marketOrderSize>0$.\footnote{
Our identification is kept in place in trading venues
where traders can actually submit market orders alongside limit orders.
In these trading venues the actual trigger of a trade might be a limit order $(t,q,p,d)$
with non-trivial price specification $0<p<\infty$ or a genuine market order.
Because of the asserted equivalence in Proposition \ref{prop.decompositionOfLimitOrder},
referring to both these triggers as ``market orders''
does not alter the evolution of the limit order book.
}

A market order does not in general execute at a single price.
The quantities $q_i$ of \eqref{eq.fillAtLevel} are the sizes of its fills,
level by level:
a buy market order for $\quantityToLiquidate$ shares
takes $q_i$ shares at the price $\nthBestAskPrice[i]_{t-}$,
and $\sum_{i\geq1}q_i = \quantityToLiquidate$.
For a sell market order the fills are mutatis mutandis the same, on the bid side.

Let $(\arrivalTimes[]_n, \marketOrderSize[n], d_n)$ be the stream of market orders
of a trading epoch,
written as in Proposition \ref{prop.decompositionOfLimitOrder},
and let $q_{n,i}$ be the fills of the $n$-th of them.

\begin{defi}
The volume transacted over the window $(t-w, t]$ is
\begin{equation}
\label{eq.transactedVolume}
\volume_{t,w} := \sum_{n \, : \, t-w < \arrivalTimes[]_n \leq t} \marketOrderSize[n] ,
\end{equation}
and the value traded over the same window is
\begin{equation}
\label{eq.tradedValue}
\tradedValue_{t,w}
:=
\sum_{n \, : \, t-w < \arrivalTimes[]_n \leq t} \, \sum_{i \geq 1} q_{n,i} \, \pi_{n,i} ,
\end{equation}
where $\pi_{n,i}$ is the price of the $i$-th fill of the $n$-th market order.
\end{defi}

\begin{defi}
The volume-weighted average price over the window $(t-w,t]$ is
\begin{equation}
\label{eq.vwap}
\vwap_{t,w} := \frac{\tradedValue_{t,w}}{\volume_{t,w}} ,
\end{equation}
defined whenever $\volume_{t,w} > 0$.
\end{defi}

\begin{remark}
Notice that $\vwap_{t,w}$ is a statistic of the tape:
it is a function of the transactions, and cannot be recovered
from the sequence of configurations alone,
since a level shrinks by cancellation as well as by execution.
The sweep cost of \eqref{eq.sweepCost}, by contrast, is a function of one configuration,
and is defined whether or not the order it prices is ever submitted.
\end{remark}
